
\documentclass[12pt,a4paper]{article}
\usepackage[utf8]{inputenc}
\usepackage[spanish]{babel}
\usepackage{amsmath}
\usepackage{amsfonts}
\usepackage{amssymb}
\usepackage{graphicx}
\usepackage{hyperref}

\title{\textbf{Cosmología Informacional Unificada}\\ \large Geometría de Nodos de Orientación Angular y la Dinámica de la Consciencia en el Deep Hilbert Space}
\author{Pablo Sebastián Martín}
\date{Abril 2026}

\begin{document}

\maketitle

\begin{abstract}
La presente investigación propone un marco teórico unificado que integra la Teoría de la Información Duodecimal ($EIT_{12}$) con la física de los Nodos de Orientación Angular (NOA). Superando el materialismo reduccionista, postulamos que el Cosmos es un continuo espacio-temporal decadimensional. A través del análisis de la dinámica en el Deep Hilbert Space (DHS), se demuestra que la consciencia opera como un factor no-lineal en la evolución sistémica, manifestado mediante el colapso de la probabilidad modificada $P'(s)$.
\end{abstract}

\section{Introducción: La Ulterioridad y el Paradigma Informacional}
La crisis actual de la física teórica exige una transición ontológica: el paso de una realidad basada en la materia a una fundamentada en la información. Lo que percibimos como "materia" es el resultado de la configuración angular de nodos informacionales.

\section{Física de los Nodos de Orientación Angular (NOA)}
El concepto de "partícula puntual" es una abstracción errónea derivada de una percepción tridimensional limitada. La unidad constituyente del Sistema Universal es el Nodo de Orientación Angular (NOA), cuya realidad se define por su configuración de fases angulares en un marco decadimensional.

\section{La Ecuación de Probabilidad Modificada y el Acoplamiento Débil}
La interacción entre la consciencia y la materia se formaliza mediante el sesgo no-lineal en el colapso de estados cuánticos:

\begin{equation}
P'(s) = \frac{P(s)(1 + \gamma_{12} C(s, I))}{Z}
\end{equation}

Donde $Z$ es la función de partición en el DHS y $\gamma_{12}$ es el coeficiente de influencia de la consciencia (Acoplamiento débil).

\section{Neurofísica: La Interfaz de Transducción}
La interfaz cuántica se localiza en las estructuras subcorticales, específicamente en el plexo coroideo ventrolateral, donde la red de átomos estables opera a una frecuencia de resonancia de $6 \cdot 12^3$ Hz.

\end{document}
